\begin{lemma}[Claim 4.1(a3)--(a4): missing-color transport]
Let \(T\) be an ETT whose terminal positive slot \(n\) is a PE with
exact chosen supporting vertex \(v_n\in V(T)\).  Put
\(j=m(v_n)\leq n\), and use the active set \(L\) and support classes
\(L(w)\) from \noderef{ExtensionBaseSupportClass}.  Let \(\sigma_n\) be
a proper coloring leaving exactly \(e_0\) uncolored and
\((T_n,D_n,\varphi_n)\)-stable.

First, if \(v\in V(T_j)\) is not the supporting vertex of any active
slot, then
\[
                    \overline\sigma_n(v)
                    =\overline\varphi_{j-1}(v).                \tag{1}
\]

Second, let \(w\in V(T_n)\), and let the nonempty class \(L(w)\) have
exact increasing enumeration
\(i_0<\cdots<i_p\), and assume this is also the exact enumeration through
round (n) of all PEs supported at (w).  Put
\[
 \mathcal S_w=\bigcup_{a=0}^{p}S_{i_a},qquad
 \gamma_0=\gamma_{i_0},\qquad \delta_p=\delta_{i_p}.
\]
Then
\[
 \overline\sigma_n(w)\cap\mathcal S_w=\{\delta_p\},            \tag{2}
\]
and
\[
 \overline\varphi_{j-1}(w)
   =(\overline\sigma_n(w)\setminus\{\delta_p\})
       \cup\{\gamma_0\}.                                     \tag{3}
\]
These are exactly Claim 4.1(a3) and (a4).
\end{lemma}
\begin{proof}
We first record the stage inclusions that will be used at each invocation
of stability.  The certificate \noderef{ExtendedTashkinovTree} contains a
\noderef{TashkinovSeries}; by that definition, every slot has a certified
\noderef{TashkinovSeriesStep}.  At a positive slot the initial constructor
is impossible, and each of the SE, PE, and RE constructors contains an
ordered-extension certificate from the current tree to the next tree.
By \noderef{TreeSequenceInitialSegment}, such a certificate preserves every
ordered vertex entry, and the definition \noderef{TreeSequenceVertexSet}
then gives inclusion of the corresponding vertex sets.  Composing these
inclusions for the slots from \(j\) through \(n-1\) gives
\[
                         V(T_j)\subseteq V(T_n).                \tag{4}
\]
If \(j=n\), this inclusion is reflexive.  Here the identification of the
series trees with \(T_j\) and \(T_n\) is exactly the ordered ladder-stage
convention of \noderef{ExtendedTashkinovLadderStage}; in the Lean indexing,
the indices called \(jSlot\) and \(q\) have values \(j-1\) and \(n-1\).
The further initial-segment field \(T_n\subset T\) in
\noderef{ExtendedTashkinovTree} also shows that every vertex of \(T_j\)
belongs to the displayed final ETT.

The definition \noderef{ETTFirstAppearanceIndex} makes \(m(x)\) the least
subscript of an ordered ladder stage containing \(x\).  Consequently, for
every \(x\in V(T_j)\), the stage \(T_j\) itself is a candidate in that
minimum and
\[
                              m(x)\leq j.                       \tag{5}
\]
Conversely, if \(x\in V(T)\) and \(m(x)\leq j\), then the least stage
containing \(x\) is one of \(T_1,\ldots,T_j\); composing the same certified
ordered extensions from that stage through \(T_j\) gives
\[
                              x\in V(T_j).                      \tag{6}
\]

The vertex clause in the definition \noderef{StableColoring}, applied on
the stable tree \(T_n\), gives
\[
                  \overline\sigma_n(x)=\overline\varphi_n(x)  \tag{7}
\]
for every \(x\in V(T_n)\).  In particular, (4) makes (7) available at
every vertex of \(T_j\).

Now let \(v\in V(T_j)\) satisfy the first hypothesis.  Suppose a round
\(s\), with \(j\leq s\leq n\), were a PE supported at \(v\).  The lower
bound is the first clause in \noderef{ExtensionBaseActiveIndices}; its
second clause is the assumed PE kind; and its existential third clause is
witnessed by \(v\in V(T)\), the exact chosen-support equality, and (5).
Thus \(s\) is active.  This contradicts the hypothesis that no active
round has chosen vertex \(v\).  Hence no PE in the full nonempty interval
from \(j\) through \(n\) is supported at \(v\).  Applying
\noderef{HkuUnsupportedVertexColors}, with starting-tree membership
\(v\in V(T_j)\), yields
\[
                  \overline\varphi_{j-1}(v)
                    =\overline\varphi_n(v).                   \tag{8}
\]
Equations (7) and (8) prove (1).

We next establish all premises needed for the two endpoint comparisons for
the class \(L(w)\).  Write \(b_a=i_a-1\) for the zero-based Lean round in
the given enumeration.  Its range equality at \(b_0\) says
\(b_0\in L(w)\).  Unfolding \noderef{ExtensionBaseSupportClass} gives both
\(b_0\in L\) and the exact equality saying that the chosen vertex at this
PE is \(w\).  Unfolding \noderef{ExtensionBaseActiveIndices} at \(b_0\)
then gives
\[
 b_0+1=i_0\geq j,\qquad w\in V(T),\qquad m(w)\leq j,            \tag{9}
\]
as well as the PE-kind equality.  In the existential clause of the active
set, its chosen vertex equals \(w\) because both optional chosen-vertex
values equal the one stored at \(b_0\).  Equation (6) applied to (9) now
supplies the required starting-tree fact
\[
                              w\in V(T_j).                     \tag{10}
\]

We prove the early endpoint equality, separating the possible empty
interval.  If \(i_0=j\), then the two endpoint colorings are the same
indexed coloring, so
\(\overline\varphi_{j-1}(w)=
  \overline\varphi_{i_0-1}(w)\) is reflexive.  Suppose instead that
\(j<i_0\).  The interval of manuscript rounds from \(j\) through
\(i_0-1\) is then nonempty.  Its zero-based indices run from \(j-1\)
through \(b_0-1=i_0-2\).  No PE in this interval is supported at \(w\):
otherwise the exact enumeration through \(n\) would identify its
zero-based index with some \(b_a\), while strict increase of the
enumeration makes \(b_0\) its least member.  This contradicts that the
index is strictly less than \(b_0\).  Apply
\noderef{HkuUnsupportedVertexColors} to this interval, whose starting
membership is (10) and whose final zero-based index is \(b_0-1\).
Its conclusion is
\[
                  \overline\varphi_{j-1}(w)
                    =\overline\varphi_{i_0-1}(w).              \tag{11}
\]

For the late endpoint, exactness of the enumeration at its last member
\(b_p=i_p-1\) gives \(b_p\leq n-1\), the PE-kind equality at the positive
slot \(i_p\), and the exact chosen vertex \(w\).  If \(i_p=n\), the desired
equality
\(\overline\varphi_{i_p}(w)=\overline\varphi_n(w)\) is reflexive.
Suppose \(i_p<n\).  Inspect the certificate at the positive slot \(i_p\),
supplied by \noderef{TashkinovSeries}.  The exact PE-kind equality rules
out the initial, SE, and RE constructors of
\noderef{TashkinovSeriesStep}, so this certificate is its PE constructor.
The constructor's exact chosen-vertex field identifies the vertex in its
\noderef{MaximumDefectiveChoice} with \(w\).  That vertex is defective in
the current tree; by the membership clause of \noderef{DefectiveVertex},
\(w\in V(T_{i_p})\).  The PE constructor's ordered-extension field, followed
by \noderef{TreeSequenceInitialSegment} and
\noderef{TreeSequenceVertexSet}, therefore gives
\[
                              w\in V(T_{i_p+1}).                \tag{12}
\]
The interval from \(i_p+1\) through \(n\) is nonempty.  It contains no PE
supported at \(w\): exactness would identify its zero-based index with
some \(b_a>b_p\), contrary to the fact that \(b_p\) is the last member of
the strictly increasing enumeration.  Thus
\noderef{HkuUnsupportedVertexColors}, starting with (12), gives
\[
                  \overline\varphi_{i_p}(w)
                    =\overline\varphi_n(w).                   \tag{13}
\]

It remains to combine the endpoint equalities with the two already
established PE identities.  Apply \noderef{SupportingPEStoredIntersection}
to the given exact enumeration and terminal bound \(n\).  It gives
\[
                  \overline\varphi_n(w)\cap\mathcal S_w
                    =\{\delta_p\}.                            \tag{14}
\]
The hypothesis \(w\in V(T_n)\) permits (7) at \(w\); substituting it into
(14) proves (2).  Next, \noderef{SupportingPEMissingTransport} applied to
the same enumeration gives
\[
 \overline\varphi_{i_0-1}(w)
   =(\overline\varphi_{i_p}(w)\setminus\{\delta_p\})
      \cup\{\gamma_0\}.                                     \tag{15}
\]
Substitute (11) on the left of (15), then (13) and (7) on its right.
The resulting equality is exactly (3), which completes the proof.
\end{proof}
